% GENERATED FILE. Edit canonical lessons, then run npm run generate:books.
% canonical-source-hash: fnv1a64:f9f7dd0764e5fb1b
% canonical-lessons: RU-C395-sberezheniya, RU-C395-vkladyvat, RU-C395-vklad, RU-C395-stoimost, RU-C395-byudzhet

\chapter{Savings and budgets}
\label{ch:ru-savings-and-budgets}
\hlchaptermodality{395}

\begin{chapteropening}
\textbf{By the end of this chapter:} \emph{I can talk about savings and a bank deposit, investing money, the cost of something and a budget.}

Complete the last lesson of chapter 395: I can talk about savings and a bank deposit, investing money, the cost of something and a budget.
\end{chapteropening}

\section[\texorpdfstring{\ru{сбережения} (sberezhéniya)}{сбережения (sberezhéniya)}]{\ru{сбережения} (sberezhéniya) — savings}

\label{lesson:RU-C395-sberezheniya}



\begin{quote}
Before the new one: say the Russian for to give back, then the Russian for a per cent.
\end{quote}

\subsection*{You'll want to know: \ru{сбережения}}
\textbf{\ru{сбережения}} (\emph{sberezhéniya}) — \textquotedblleft{}savings\textquotedblright{}.

\begin{quote}
\ru{Все} \ru{сбережения} \ru{он} \ru{потратил} \ru{на} \ru{учёбу}.
\end{quote}

\begin{quote}
\emph{Vse sberezhéniya on potrátil na uchyóbu.}
\end{quote}

\begin{quote}
He spent all his savings on his studies.
\end{quote}

Plural. From \emph{\ru{сберегать}}, to save; \emph{\ru{беречь}} has the same root.

\begin{cousinweb}
\textbf{\ru{сбережения}} is built on \emph{\ru{сберечь}} \textquotedblleft{}to save, to keep\textquotedblright{}, from \emph{\ru{беречь}} \textquotedblleft{}to protect\textquotedblright{}, inherited from Proto-Slavic \emph{*berťi} (EtymDB).
\end{cousinweb}

\begin{grammarlens}[title={What you've built}]
One more word for telling, explaining and describing at B1.
\end{grammarlens}

\subsection*{Your turn}
\begin{itemize}
\raggedright
  \item \emph{Say it:} \emph{sberezhéniya}
  \item \emph{Say it:} \emph{sberezhéniya}, once more
  \item \emph{Say it:} a sentence of your own with \emph{sberezhéniya}
  \item \emph{Recall:} say the Russian for to give back, then the Russian for a per cent, then say \emph{sberezhéniya} again
\end{itemize}

\begin{tcolorbox}[breakable,colback=teal!4,colframe=teal!35!black,title={Before you move on}]
What does \ru{сбережения} mean? (\textquotedblleft{}Savings\textquotedblright{}.) Say it once more.
\end{tcolorbox}
\section[\texorpdfstring{\ru{вкладывать} (vkládyvat)}{вкладывать (vkládyvat)}]{\ru{вкладывать} (vkládyvat) — to invest, to put in}

\label{lesson:RU-C395-vkladyvat}



\begin{quote}
Before the new one: say the Russian for a per cent, then the Russian for savings.
\end{quote}

\subsection*{You'll want to know: \ru{вкладывать}}
\textbf{\ru{вкладывать}} (\emph{vkládyvat}) — \textquotedblleft{}to invest, to put in\textquotedblright{}.

\begin{quote}
\ru{Он} \ru{вкладывает} \ru{деньги} \ru{в} \ru{недвижимость}.
\end{quote}

\begin{quote}
\emph{On vkládyvayet déngi v nedvízhimost.}
\end{quote}

\begin{quote}
He invests money in property.
\end{quote}

Imperfective; perfective \emph{\ru{вложить}}. \emph{\ru{Вкладывать} \ru{в}} + accusative.

\begin{cousinweb}
\textbf{\ru{вкладывать}} is built from \emph{\ru{в}-} \textquotedblleft{}in\textquotedblright{} and the root of \emph{\ru{класть}} \textquotedblleft{}to lay\textquotedblright{}, inherited from Proto-Slavic \emph{*klasti}: to lay money into something.
\end{cousinweb}

\begin{grammarlens}[title={What you've built}]
One more word for telling, explaining and describing at B1.
\end{grammarlens}

\subsection*{Your turn}
\begin{itemize}
\raggedright
  \item \emph{Say it:} \emph{vkládyvat}
  \item \emph{Say it:} \emph{vkládyvat}, once more
  \item \emph{Say it:} a sentence of your own with \emph{vkládyvat}
  \item \emph{Recall:} say the Russian for a per cent, then the Russian for savings, then say \emph{vkládyvat} again
\end{itemize}

\begin{tcolorbox}[breakable,colback=teal!4,colframe=teal!35!black,title={Before you move on}]
What does \ru{вкладывать} mean? (\textquotedblleft{}To invest, to put in\textquotedblright{}.) Say it once more.
\end{tcolorbox}
\section[\texorpdfstring{\ru{вклад} (vklad)}{вклад (vklad)}]{\ru{вклад} (vklad) — a deposit; a contribution}

\label{lesson:RU-C395-vklad}



\begin{quote}
Before the new one: say the Russian for savings, then the Russian for to invest.
\end{quote}

\subsection*{You'll want to know: \ru{вклад}}
\textbf{\ru{вклад}} (\emph{vklad}) — \textquotedblleft{}a deposit; a contribution\textquotedblright{}.

\begin{quote}
\ru{Я} \ru{открыл} \ru{вклад} \ru{в} \ru{банке}.
\end{quote}

\begin{quote}
\emph{Ya otkrýl vklad v bánke.}
\end{quote}

\begin{quote}
I opened a deposit account at the bank.
\end{quote}

Masculine. \emph{\ru{Вклад} \ru{в}} + accusative — a contribution to: \emph{\ru{большой} \ru{вклад} \ru{в} \ru{науку}}.

\begin{cousinweb}
\textbf{\ru{вклад}} is formed from \emph{\ru{вкладывать}} \textquotedblleft{}to put in\textquotedblright{}, from \emph{\ru{в}-} and \emph{\ru{класть}} \textquotedblleft{}to lay\textquotedblright{}, Proto-Slavic \emph{*klasti}.
\end{cousinweb}

\begin{grammarlens}[title={What you've built}]
One more word for telling, explaining and describing at B1.
\end{grammarlens}

\subsection*{Your turn}
\begin{itemize}
\raggedright
  \item \emph{Say it:} \emph{vklad}
  \item \emph{Say it:} \emph{vklad}, once more
  \item \emph{Say it:} a sentence of your own with \emph{vklad}
  \item \emph{Recall:} say the Russian for savings, then the Russian for to invest, then say \emph{vklad} again
\end{itemize}

\begin{tcolorbox}[breakable,colback=teal!4,colframe=teal!35!black,title={Before you move on}]
What does \ru{вклад} mean? (\textquotedblleft{}A deposit; a contribution\textquotedblright{}.) Say it once more.
\end{tcolorbox}
\section[\texorpdfstring{\ru{стоимость} (stóimost)}{стоимость (stóimost)}]{\ru{стоимость} (stóimost) — cost, value}

\label{lesson:RU-C395-stoimost}



\begin{quote}
Before the new one: say the Russian for to invest, then the Russian for a deposit.
\end{quote}

\subsection*{You'll want to know: \ru{стоимость}}
\textbf{\ru{стоимость}} (\emph{stóimost}) — \textquotedblleft{}cost, value\textquotedblright{}.

\begin{quote}
\ru{Стоимость} \ru{билета} — \ru{тысяча} \ru{рублей}.
\end{quote}

\begin{quote}
\emph{Stóimost biléta — týsyacha rubléy.}
\end{quote}

\begin{quote}
The ticket costs a thousand roubles.
\end{quote}

Feminine. From \emph{\ru{стоить}}, to cost.

\begin{cousinweb}
\textbf{\ru{стоимость}} is built on \emph{\ru{стоить}} \textquotedblleft{}to cost\textquotedblright{}, which perhaps came through Polish and copies the sense of Latin \emph{cōnstāre} \textquotedblleft{}to stand firm; to cost\textquotedblright{}, by a usual account; it goes with \emph{\ru{стоять}} \textquotedblleft{}to stand\textquotedblright{}.
\end{cousinweb}

\begin{grammarlens}[title={What you've built}]
One more word for telling, explaining and describing at B1.
\end{grammarlens}

\subsection*{Your turn}
\begin{itemize}
\raggedright
  \item \emph{Say it:} \emph{stóimost}
  \item \emph{Say it:} \emph{stóimost}, once more
  \item \emph{Say it:} a sentence of your own with \emph{stóimost}
  \item \emph{Recall:} say the Russian for to invest, then the Russian for a deposit, then say \emph{stóimost} again
\end{itemize}

\begin{tcolorbox}[breakable,colback=teal!4,colframe=teal!35!black,title={Before you move on}]
What does \ru{стоимость} mean? (\textquotedblleft{}Cost, value\textquotedblright{}.) Say it once more.
\end{tcolorbox}
\section[\texorpdfstring{\ru{бюджет} (byudzhét)}{бюджет (byudzhét)}]{\ru{бюджет} (byudzhét) — a budget}

\label{lesson:RU-C395-byudzhet}



\begin{quote}
Before the new one: say the Russian for a deposit, then the Russian for cost.
\end{quote}

\subsection*{You'll want to know: \ru{бюджет}}
\textbf{\ru{бюджет}} (\emph{byudzhét}) — \textquotedblleft{}a budget\textquotedblright{}.

\begin{quote}
\ru{У} \ru{нас} \ru{маленький} \ru{семейный} \ru{бюджет}.
\end{quote}

\begin{quote}
\emph{U nas málenkiy seméynyy byudzhét.}
\end{quote}

\begin{quote}
We have a small family budget.
\end{quote}

Masculine. \emph{\ru{Государственный} \ru{бюджет}} — the state budget.

\begin{cousinweb}
\textbf{\ru{бюджет}} is usually traced to English \emph{budget}, from Old French \emph{bougette} \textquotedblleft{}a little bag\textquotedblright{}, from Latin \emph{bulga} \textquotedblleft{}a leather bag\textquotedblright{} (Lewis \& Short): first the bag the money was kept in.
\end{cousinweb}

\begin{grammarlens}[title={What you've built}]
One more word for telling, explaining and describing at B1.
\end{grammarlens}

\subsection*{Your turn}
\begin{itemize}
\raggedright
  \item \emph{Say it:} \emph{byudzhét}
  \item \emph{Say it:} \emph{byudzhét}, once more
  \item \emph{Say it:} a sentence of your own with \emph{byudzhét}
  \item \emph{Recall:} say the Russian for a deposit, then the Russian for cost, then say \emph{byudzhét} again
\end{itemize}

\begin{tcolorbox}[breakable,colback=teal!4,colframe=teal!35!black,title={Before you move on}]
What does \ru{бюджет} mean? (\textquotedblleft{}A budget\textquotedblright{}.) Say it once more.
\end{tcolorbox}
